Identifying maternal drug use during pregnancy enables more appropriate medical care and linkage to support services. However, drug tests may also be used for non-therapeutic purposes, such as criminal justice and child protective services involvement. Racial disparities in drug testing could undermine patient trust in the health care system, negatively impacting health outcomes. We used electronic medical records from $N = 140,562$ infants born between 2019-2023 in Indiana hospitals to measure meconium and umbilical cord tests for in-utero exposure to opioids and other illegal drugs. Black infants were tested at higher rates (16\%) compared to White infants (7\%), but both groups had a 20\% positive testing rate. We used inverse propensity score weights to decompose the testing gap into a portion explained by measured risk factors and an unexplained portion. Measured risk differences account for approximately 38\% of the testing gap. The unexplained portion of the gap suggests that physicians use a different threshold for testing White and Black infants, or that physician decisions are motivated by additional risk factors that are not included in our adjustment.

\bigskip
